
\section{Observable Projection}

The observable distance perturbation is not identified directly with the scalar amplitude. At linear order it is a light-cone projection of metric, velocity, and scalar perturbations:
\[
\frac{\delta D_L}{\bar D_L}
=
\int d\chi\,
\left[
W_\Phi\Phi+
W_\Psi\Psi+
W_vv_\parallel+\cdots
\right].
\]
If the selected \(n=2\) curvature-scale sector admits a physical gauge-invariant metric response, its exact scalar angular dependence can factor as
\[
\hat{\mathbf n}\cdot\hat{\mathbf p},
\]
leading schematically to
\[
\frac{\delta D_L}{\bar D_L}
=
\epsilon_D(z)\,
\hat{\mathbf n}\cdot\hat{\mathbf p}.
\]
The existence and normalization of \(\epsilon_D(z)\) are therefore bridge outputs to be derived, not consequences of the scalar harmonic geometry alone.

For supernovae,
\[
\mu=5\log_{10}D_L+\mathrm{const},
\]
and therefore a small magnitude dipole \(A_\mu\) corresponds to
\[
\frac{\delta D_L}{D_L}
=
\frac{\ln10}{5}A_\mu\cos\theta.
\]

For structure growth we define
\[
k^2\Psi
=
-4\pi Ga^2\mu(k,a)\rho_m\delta_m,
\qquad
\Phi=\gamma(k,a)\Psi,
\]
and
\[
\Sigma(k,a)
=
\frac{\mu(k,a)}2[1+\gamma(k,a)].
\]
The matter growth factor obeys
\[
D''+
\left(2+\frac{H'}H\right)D'
-\frac32\Omega_m(a)\mu(k,a)D=0,
\]
with prime denoting \(d/d\ln a\). The observable is
\[
f\sigma_8(k,z)
=
\frac{D'}{D}D\,\sigma_{8,\rm ini}
=
D'(k,z)\sigma_{8,\rm ini}.
\]

The quasi-static limit is used only on scales inside the scalar sound horizon \citep{SawickiBellini2015}. Horizon-scale \(n=2\) observables require the full constrained long-wavelength response and must not be obtained by extrapolating quasi-static formulas into the first physical curvature-scale mode.

% BEGIN R1_COSMOLOGY_CLOSURE_V15
\subsection{Current action-native observable status}
\label{sec:r1-observable-status-v15}

The present R1 observable construction is evaluated on the canonical
two-component state
\[
X_2=(q_2,\Pi_2),
\qquad
q_2=\zeta_{\rm anchor},
\qquad
\Pi_2=2a_{\rm anchor}^3\mathcal G_{S,2}\dot\zeta_{\rm anchor}.
\]
Legacy response rows expressed in $(\zeta,d\zeta/dN)$ coordinates are not
interchangeable with these canonical rows without the explicit basis
transformation.

For standard candles the action-native fixed-observed-redshift luminosity-distance
row is denoted $\mathcal F_{D_L}(z)$. Distance duality then closes the
transverse late-time BAO distance response,
\[
\boxed{\mathcal F_{D_M/r_d}(z)=\mathcal F_{D_L}(z)},
\qquad d\ln r_d=0.
\]
This statement concerns the observed distance to the standard ruler, not an
intrinsic deformation of the sound horizon.

The historical closed-background expansion
\[
\frac{\Delta r_{\rm BAO}}{r_{\rm BAO}}
\simeq -\frac{K\chi^2(z)}{6}
\]
is therefore retained only as a background distance-projection expression. For
the frozen R1 radius $R_H=33.153\,\mathrm{Gpc}$, curvature across a
$150\,\mathrm{Mpc}$ ruler itself would produce only
$-3.412e-06$ at leading order, whereas the Gpc-baseline projection
at $z=0.93$ is $-1.562e-03$.

The radial observable remains an explicit open kernel:
\[
\mathcal F_{D_H/r_d}(z)=-\mathcal F_H(z)
\]
on the same late-time branch. We do not promote a numerical derivative of the
transverse distance kernel to $\mathcal F_H$, because the transverse light-cone
response contains metric, velocity, and integrated optical terms that need not
reduce to a single radial remapping. Survey window and covariance projection
are downstream of these cosmic kernels.
% END R1_COSMOLOGY_CLOSURE_V15
